% Self-contained fragment for the draft of article II, September 8, 2026.
% Provenance: RESULTADO_RECUPERADO and FORMALIZACION_REUNIDA from the owners
% B07/c41_ckm.tex and C02/74_mezcla_contenido_exclusivo_post_rev2.tex in
% BORRADOR_02; integral 2249, chaps. 49, 92.6, 92.24, and 93.10--93.11.
% The originals and their complete proof support are neither replaced nor edited.
% Bibliographic keys used: ckm, mezcla, pmnsespectral, cpt.
% Requires: amsmath, amssymb, amsthm; theorem/proposition/definition/remark.

\section{CKM mixing angles, CP phase, and leptonic transport}
\label{sec:ii-propagacion-angular}

\subsection{Common antecedent and separation of the two transports}

The article's common foundation produces the enriched state through
the two APP evaluations, the oriented selection of TRIT, and the
effective composition of selection, transport, carry, memory, and return
in TPK. Angular representations do not replace that state: they are readers
of its record. In particular, the prolongation
\(w_6\to w_{12}\to w_{18}\to w_{24}\to w_{30}\to R_{36}\to G_9\)
preserves the provenance of the numerical readings, and the holonomy
\(\Gamma_9\) acts before the reduced representation \(K_{\rm ph}\).
Phase return does not restart the record.

Two classes of objects are received here from that same antecedent. The first
contains the direct angle, the counter-angle, and the torsional defect
\(\Delta_4\), already constructed through their own readers. The second contains
the surviving extensions, their observable routes, the dodecaphasic
record, and the operators of the material fibers. The
extension--route--record--signature--character--towers chain preserves the atlas
\(81/56/13/324\) and distinguishes the null chart from the positive branch.
The mixing angles do not retrospectively choose those routes or generate
the mass eigenvalues.

Quark transport uses a rational angular chart and its complex
unitary realization. Leptonic transport uses the spectral frames
of the charged and neutral fibers. Both come from the common discrete
structure, but their domains, readers, and changes of basis are distinct.
The Barbero--Immirzi functional and mixing are not defined by one
another either: they receive typed representations of the angular and incidence antecedent.
Their relationship is expressed through those maps, not through a
retrospective formula using a Barbero value as a mixing selector.

\begin{remark}[Sectorial rules and evaluation of the reader]
The three sectorial rules and their evaluation are developed below
in Section~\ref{sec:ii09-construccion-sectorial}. Theorem
\ref{thm:ii09-matriz-sectorial} obtains the columns as images of the
angular basis. Subsequent proofs establish inversion of the chart,
sheet transport, the metric, and the complex realization.
The choice of rules and evaluation of their coefficients are two
distinct mathematical questions: the starting system is declared in
Definition~\ref{def:ii09-sistema-sectorial}, and the theorem's proof
uses that complete system. The preceding composition that selects those
rules from incidence requires its own proof.
\end{remark}

% Separate contribution, 9 September 2026.
% RESULTADO_RECUPERADO: sectorial construction from source 08_ckm_cp.tex.
% FORMALIZACION_REUNIDA: parity decomposition and representational compatibility.
% Insert before the rational CKM chart; retain the subsequent development.
% Provenance and exact scope are in 08b_procedencia_angular.md.

\subsection{Sectorial incidence and construction of the angular reader}
\label{sec:ii09-construccion-sectorial}

The angular construction uses the action and closure representations
previously obtained from APP--TRIT--TPK. The direct angle,
counter-angle, and torsional defect preserve their own domains.
The quark reader combines these three coordinates through sectorial
incidence; its evaluation precedes complex amplitudes and
metrological comparison. Hexadic--octadic inclusion enters
the third sector through the difference of its cardinalities.

% ARQUITECTURA_AUTORAL_PREEXISTENTE: bidirectional half-square of N42.
% FORMALIZACION_NUEVA: pair groupoid and stabilizer weighting.
\subsubsection{Marked incidence and bidirectional normalization}
\label{sec:ii09-normalizacion-bidireccional}

The reader invariants preserve the nature of the object from which
they arise. The marked intersection $B_K=H_e\cap P_\pi$ has cardinality
$b=4$; the integer $q=7$ is the coordinate of the pivot $p_\varphi$ in
the oriented numbering. The determination order of the Steiner
design is $p=5$. The hexadic and octadic supports have cardinalities $h_6=6$
and $o_8=8$, and the tritic orientation alphabet has cardinality $t=3$.
A support cardinality and a pivot coordinate are used through
their respective maps, even when both are expressed as integers.

Normalization of the pentagonal self-crossing admits the following
explicit realization. Let $P$ be the five-position support of the closure,
and let $X=P\times P$ be its space of ordered pairs. Bidirectional
inversion on this space is
\[
 \jmath:X\longrightarrow X,\qquad \jmath(x,y)=(y,x),
 \qquad\jmath^2=I.
\]
The stabilizer information of each orbit is preserved. For an
orbit $[z]$, define its weight as
$1/|\operatorname{Stab}_{C_2}(z)|$. The sum of these weights is the
cardinality of the action groupoid $X/\!/C_2$.

\begin{proposition}[Orbital measure of bidirectional self-crossing]
\label{prop:ii09-semicuadrado-orbital}
For a finite support of cardinality $p$, inversion of ordered pairs
determines
\[
 |(P\times P)/\!/C_2|
 =\sum_{[z]\in(P\times P)/C_2}
      \frac1{|\operatorname{Stab}_{C_2}(z)|}
 =\frac{p^2}{2}.
\]
For $p=5$, the self-crossing has ten free orbits and five fixed
orbits, and its orbital measure is $10+5/2=25/2$.
\end{proposition}
\begin{proof}
Each diagonal pair $(x,x)$ forms a fixed orbit, whose stabilizer
has order two; there are $p$ such orbits. The remaining $p(p-1)$ pairs
are grouped in twos, with trivial stabilizer. Their contribution is
therefore $p(p-1)/2$, and that of the diagonal is $p/2$. The sum is
$p^2/2$. Equivalently, the orbit--stabilizer identity gives
$1/|\operatorname{Stab}(z)|=|[z]|/2$; summing over all orbits
produces $|P\times P|/2$.
\end{proof}

The measure explicitly preserves the contribution of fixed points.
The ordinary orbit set has cardinality $p(p+1)/2$, equal to
fifteen in the pentagonal closure. The orbital measure and that cardinality
are two determinate readings of the same inversion. The reader's
bidirectional normalization uses the former: the ten free orbits
have unit weight, and the five diagonals have half-unit weight.

This realization specifies the factor $p^2/2$ used in the following
sectorial system. The assignment of that contribution to sector $23$,
its negative orientation, and the other two incidence rules remain
specified by the complete sectorial system. The proof establishes the
normalization of self-crossing and preserves the selection of the
angular functionals as a distinct matter.


% FORMALIZACION_NUEVA of the lift of N39, compatible with (rho9,q9),
% T_d and Upd_t; preexisting remainder/quotient and orientation architecture.
% Separate research fragment. Does not select the CKM N42 rules.
\subsubsection{Carry transport and integer holonomy of the arithmetic crossing}
\label{subsec:transporte-entero-cruce}

The modular reading of the crossing admits an integral realization retaining
the information accumulated while crossing the edge. This realization
is developed here through transport quotients. Its domain is a lifted
chart of the APP support; the subsequent assignment of a region to an angular
sector retains its own incidence map.

\paragraph{Nonadic chart and lifted transport.}
We use the exact decomposition
\[
 \rho_9(n)=1+((n-1)\bmod9),\qquad
 q_9(n)=\left\lfloor\frac{n-1}{9}\right\rfloor,
 \qquad n=\rho_9(n)+9q_9(n).
\]
A lifted position is
$z=(x,y,q_x,q_y)\in D_9^2\times\mathbb Z^2$, where
$X=x+9q_x$ and $Y=y+9q_y$. The integer increment $u$ of the first
coordinate acts by
\[
 T^x_u(x,y,q_x,q_y)=
 \bigl(\rho_9(x+u),y,q_x+c_9(x;u),q_y\bigr),\qquad
 c_9(x;u)=\left\lfloor\frac{x+u-1}{9}\right\rfloor.
\]
The second coordinate is transported in the same way. Cardinal
advances correspond to $u=\pm1$; the zero increment preserves the state.
The cocycle satisfies
\[
 c_9(x;u+v)=c_9(x;u)+c_9(\rho_9(x+u);v).
\]
Indeed, both sides are the quotient of the unique decomposition of
$x+u+v$ with remainder in $D_9$. Consequently,
$T^x_vT^x_u=T^x_{u+v}$ and $T^x_{-u}=(T^x_u)^{-1}$; the transports
of the two coordinates commute. The state preserves the number of edge
crossings even when the visible position returns.

\paragraph{Lifted sheets and update of records.}
The integral realization is fixed by prolonging the APP arithmetic
operations to the transported coordinates:
\begin{align*}
 \widetilde S(z)&=X+Y=x+y+9(q_x+q_y),\\
 \widetilde P(z)&=XY
 =xy+9(q_xy+q_yx)+81q_xq_y.
\end{align*}
At each position the record of both sheets is preserved:
\[
 \bigl(\rho_9(\widetilde S),q_9(\widetilde S),
       \rho_9(\widetilde P),q_9(\widetilde P)\bigr).
\]
If $e:z\to z'$ is a link and $T$ the selected sheet, its full increment
is obtained from the record before and after transport:
\begin{equation}
 \delta_e\widetilde T=
 \rho_9(\widetilde T(z'))-\rho_9(\widetilde T(z))
 +9\bigl(q_9(\widetilde T(z'))-q_9(\widetilde T(z))\bigr).
 \label{eq:cruce-incremento-registro}
\end{equation}
This identity makes preservation explicit: the residue jump and
quotient increment together reconstruct the variation of the sheet.
Sheet selection, cardinal transport, and record
update are distinct operations. In this chart, the ordered sheet
selection $(T_x,T_y)$ is first composed with $T_d$ and the update
of its quotients; the formula does not postulate that the entire TPK is reduced
to this local representation.

\begin{proposition}[Integral curvature and boundary transport]
\label{prop:cruce-stokes-integral}
Define
$\widetilde A_x=\delta_x\widetilde T_x$ and
$\widetilde A_y=\delta_y\widetilde T_y$ through
\eqref{eq:cruce-incremento-registro}, and let
\[
 \widetilde F_{xy}(X,Y)=
 \widetilde A_x(X,Y)+\widetilde A_y(X+1,Y)
 -\widetilde A_x(X,Y+1)-\widetilde A_y(X,Y).
\]
Homogeneous choices have zero curvature. The mixed choices
$(\widetilde S,\widetilde P)$ and
$(\widetilde P,\widetilde S)$ have curvature
$+1$ and $-1$, respectively, as integers. For the positive boundary of a rectangle
$R_{a,b}$ with integer sides $a,b>0$,
\[
 \widetilde W_{SP}(\partial R_{a,b})=ab,\qquad
 \widetilde W_{PS}(\partial R_{a,b})=-ab.
\]
The result is independent of the base point and the number of crossings
of nonadic seams.
\end{proposition}

\begin{proof}
Exact reconstruction of the two sheets gives
\[
 \delta_x\widetilde S=\delta_y\widetilde S=1,
 \qquad\delta_x\widetilde P=Y,
 \qquad\delta_y\widetilde P=X.
\]
Thus the $SP$ connection is $(1,X)$, with curvature
$1+(X+1)-1-X=1$; the $PS$ connection is $(Y,1)$, with curvature
$Y+1-(Y+1)-1=-1$. For a single sheet, mixed
differences cancel. Summing the curvatures of the $ab$ plaquettes cancels each
interior link with its inverse and preserves the boundary links.
This yields $\pm ab$ in $\mathbb Z$. The transported quotients
ensure the same identity holds across any seam.
\end{proof}

The boundary accumulator can in turn be retained in nonadic form:
$W=r_W+9q_W$. For an integer contribution $w_e$, it is updated by
\[
 r_W'=\rho_9(r_W+w_e),\qquad
 q_W'=q_W+\left\lfloor\frac{r_W+w_e-1}{9}\right\rfloor.
\]
Composition by links provides the same sum as direct integer
evaluation. In this chart, integer zero has coordinates
$(r_W,q_W)=(9,-1)$; this is an arithmetic representation of the accumulator,
without further identification with the distinct oriented zeros of TRIT.

\paragraph{Three operations on orientation.}
Let $\gamma$ be a closed cardinal route. Its time reversal
$\gamma^{-1}$ reverses the order of links and the orientation of each;
by additivity,
\[
 \widetilde W_{SP}(\gamma^{-1})=-\widetilde W_{SP}(\gamma).
\]
Exchange of sheets likewise satisfies
\[
 \widetilde W_{PS}(\gamma)=-\widetilde W_{SP}(\gamma).
\]
To check this on any route, note that the sum of both
connections is the exact increment of
$\widetilde S+\widetilde P=X+Y+XY$; its integral over a loop is zero.

The simultaneous half-turn of both directions, on the other hand, acts
through $J_c(X,Y)=(2c_x-X,2c_y-Y)$, with $2c\in\mathbb Z^2$,
so that it preserves the lattice chart. For a loop,
\[
 \widetilde W_{SP}(J_c\gamma)=\widetilde W_{SP}(\gamma).
\]
Indeed, the closed horizontal contribution is zero, and the vertical one
transforms as
$\sum (2c_x-X)(-\delta Y)=\sum X\delta Y$, since
$\sum\delta Y=0$. A spatial half-turn therefore preserves the
area orientation of this chart. Its combination with time reversal
or sheet exchange is computed by composing the corresponding
operations.

The difference matters for the electronic record: simultaneous reversal
of both displacements every $27$ advances produces, in
windows of $9$, the sequence
\[
 \tau_j=(-1)^{\lfloor j/3\rfloor},\qquad 0\leq j<12,
 \qquad (+,+,+,-,-,-,+,+,+,-,-,-).
\]
This sequence records advance orientations. Holonomy additionally requires
the links traversed, selected sheets, and accumulated
quotients. Transport of electronic orientation is thus preserved
without identifying a half-turn of both axes with time reversal
of a loop.

\paragraph{Application to the marked crossing and sector scope.}
For the rectangular crossing with sides $4$ and $7$, this realization evaluates
exactly
\[
 \widetilde W_{SP}=28=1+9\cdot3,\qquad
 \widetilde W_{SP}(\gamma^{-1})=-28=8+9(-4).
\]
The integer preserves the oriented area; the isolated remainder represents only
its nonadic class. For example, from $(7,1)$ the periodic
representatives of the links sum to $-35$, whereas the carry correction
adds $63$ and recovers $28$.

In the marked incidence of the angular reader, $4$ is
$|H_e\cap P_\pi|$ and $7$ is the marked pivot $p_\varphi$.
The rectangular reading supplies an integer realization of the product
$|H_e\cap P_\pi|\,p_\varphi$ once that crossing has been declared. The
map identifying that region with sector $12$ must preserve
the mark and orientation of the sector system in
Definition~\ref{def:ii09-sistema-sectorial}. Sector assignments
$12,23,13$ and their combinations of $A,h,\Delta^\circ$
are an additional specification beyond the rectangular-holonomy law.
The present result makes the area lift, carry, and
their involutions explicit, with these starting data declared.


\begin{definition}[Sectorial system of angular incidence]
\label{def:ii09-sistema-sectorial}
The sectorial system consists of the ordered sextuple of invariants
\[
 (b,p,q,h_6,o_8,t)=(4,5,7,6,8,3)
\]
and the linear maps, on angular coordinates
\((a,u,d)\),
\begin{align}
 \mathscr A_{12}^{\rm ang}(a,u,d)&=2a-pu+bq\,d,\label{eq:ii09-sector12}\\
 \mathscr A_{23}^{\rm ang}(a,u,d)&=qu-\frac{p^2}{2}\,d,\label{eq:ii09-sector23}\\
 \mathscr A_{13}^{\rm ang}(a,u,d)&=a-pb\,u-\frac{o_8-h_6}{t}\,d.
 \label{eq:ii09-sector13}
\end{align}
The symbols \(h_6,o_8\) preserve the cardinalities of the marked
incidence; \(t=3\) is the declared tritic normalization of that sector.
The sextuple and the three rules jointly form the starting
system of this construction.
\end{definition}

\begin{theorem}[Exact evaluation of the sectorial system]
\label{thm:ii09-matriz-sectorial}
The joint map
\((\mathscr A_{12}^{\rm ang},\mathscr A_{23}^{\rm ang},
\mathscr A_{13}^{\rm ang})^{\mathsf T}\) of
Definition~\ref{def:ii09-sistema-sectorial} has the unique matrix
\[
 \mathsf M_{\rm CKM}=
 \begin{pmatrix}
 2&-5&28\\
 0&7&-25/2\\
 1&-20&-2/3
 \end{pmatrix}.
\]
The three columns are the images of the coordinate vectors
\((1,0,0)\), \((0,1,0)\), and \((0,0,1)\) under these maps.
\end{theorem}
\begin{proof}
The sectorial operations give
\[
 bq=4\cdot7=28,\qquad p^2/2=5^2/2=25/2,
 \qquad pb=5\cdot4=20,\qquad
 (o_8-h_6)/t=(8-6)/3=2/3.
\]
Thus the successive images of the coordinate basis are
\[
 c_A=(2,0,1)^{\mathsf T},\qquad
 c_h=(-5,7,-20)^{\mathsf T},\qquad
 c_\Delta=(28,-25/2,-2/3)^{\mathsf T}.
\]
A linear map is determined by its images on a
basis. The matrix with these images as its columns is the one indicated.
The proof uses the complete sectorial system of the definition.
\end{proof}

On the previously produced coordinates, take
\[
 (a,u,d)=\left(A_{\deg},\frac{C^*_{\deg}}6,
                       \frac{180}{\pi}\Delta_4\right).
\]
The direct, conjugate, and torsional contributions to each angle
are thereby individualized. In particular, the difference
\(o_8-h_6=2\) is preserved in the coefficient of the torsional
defect in sector \(13\). This concrete presence makes it possible to follow
incidence through to the angular reader, keeping the sectorial
system, its evaluation, and the unitary parametrization separate.


\subsection{Rational CKM chart and angular units}

We write \(A_{\rm deg}\) and \(C^*_{\rm deg}\) for the coordinates in
degrees of the HMT angles already generated. We use the abbreviations
\begin{equation}
 h_\theta=\frac{C^*_{\rm deg}}6,
 \qquad
 \Delta_{\rm deg}=\frac{180}{\pi_{\rm HMT}}\Delta_4,
 \qquad
 x_\theta=(A_{\rm deg},h_\theta,\Delta_{\rm deg})^{\mathsf T}.
 \label{eq:ii-angular-chart}
\end{equation}
Here, \(h_\theta\) is an angular coordinate, distinct from the Planck
constant \(h\). Conversion to degrees is a subsequent change of coordinates; it
does not introduce a conventional value of \(\pi\) into the generator.

\begin{definition}[Quark angular reader]
The inherited reader \(\mathsf M_{\rm CKM}:\mathbb R^3\to\mathbb R^3\)
evaluates the coordinates
\begin{equation}
 \boldsymbol\theta_{\rm deg}=\mathsf M_{\rm CKM}x_\theta,
 \qquad
 \mathsf M_{\rm CKM}=
 \begin{pmatrix}
 2&-5&28\\
 0&7&-25/2\\
 1&-20&-2/3
 \end{pmatrix},
 \qquad
 \delta_{\rm deg}=9A_{\rm deg}.
 \label{eq:ii-ckm-map}
\end{equation}
The component order is \((12,23,13)\). Subsequent trigonometric
functions are evaluated at
\(\vartheta_{ab}=(\pi_{\rm HMT}/180)\theta_{ab,\rm deg}\) and
\(\delta=(\pi_{\rm HMT}/180)\delta_{\rm deg}\).
\end{definition}

\begin{proposition}[Rational reconstruction of the coordinates]
The reader has determinant \(-3857/6\) and inverse
\begin{equation}
 \mathsf M_{\rm CKM}^{-1}=
 \begin{pmatrix}
 1528/3857&3380/3857&801/3857\\
 75/3857&176/3857&-150/3857\\
 6/551&-30/551&-12/551
 \end{pmatrix}.
 \label{eq:ii-ckm-inverse}
\end{equation}
In particular,
\begin{align}
 A_{\rm deg}
 &=\frac{1528\theta_{12,\rm deg}+3380\theta_{23,\rm deg}
          +801\theta_{13,\rm deg}}{3857},\nonumber\\
 C^*_{\rm deg}
 &=\frac{6(75\theta_{12,\rm deg}+176\theta_{23,\rm deg}
          -150\theta_{13,\rm deg})}{3857},\nonumber\\
 \Delta_{\rm deg}
 &=\frac{6\theta_{12,\rm deg}-30\theta_{23,\rm deg}
          -12\theta_{13,\rm deg}}{551}.
 \label{eq:ii-ckm-recover}
\end{align}
The phase satisfies the exact linear check
\begin{equation}
 3857\delta_{\rm deg}
 =9(1528\theta_{12,\rm deg}+3380\theta_{23,\rm deg}
      +801\theta_{13,\rm deg}).
 \label{eq:ii-ckm-phase-check}
\end{equation}
\end{proposition}

\begin{proof}
Expansion of the determinant in \eqref{eq:ii-ckm-map} gives
\(-3857/6\). Multiplying the matrix given in
\eqref{eq:ii-ckm-inverse} by \(\mathsf M_{\rm CKM}\), in either
order, gives \(I_3\). Its rows recover the three components of
\(x_\theta\); the second is multiplied by six to recover
\(C^*_{\rm deg}\). Finally, we use
\(\delta_{\rm deg}=9A_{\rm deg}\).
\end{proof}

Inversion reconstructs the reader's coordinates once the reader is fixed. It does not
turn observed angles into antecedents of \(A\), \(C^*\), or
\(\Delta_4\).

\subsection{Geometric decomposition under sheet inversion}
\label{sec:ii09-descomposicion-paridad}

In the angular chart, sheet inversion preserves \(a,d\) and
reverses \(u\). The expressed vector admits the decomposition
\[
 \boldsymbol\theta_+=a c_A+u c_h+d c_\Delta,
 \qquad
 \boldsymbol\theta_-=a c_A-u c_h+d c_\Delta.
\]
\begin{proposition}[Even and odd components of angular transport]
\label{prop:ii09-paridad-angular}
The two vectors satisfy
\[
 \frac{\boldsymbol\theta_++\boldsymbol\theta_-}{2}
       =a c_A+d c_\Delta,\qquad
 \frac{\boldsymbol\theta_+-\boldsymbol\theta_-}{2}=u c_h.
\]
The contribution that changes sign occupies the line
\(\operatorname{span}(c_h)\); the preserved component occupies
\(\operatorname{span}(c_A,c_\Delta)\).
\end{proposition}
\begin{proof}
The sum cancels the terms \(u c_h\), and the difference cancels
the other two. The columns are linearly independent, since
\(\det\mathsf M_{\rm CKM}=-3857/6\). The two components therefore
form a direct decomposition.
\end{proof}

The same inversion acts on the Barbero--Immirzi evaluation:
it exchanges \(q_+\) and \(q_-\), and reverses the component bilinear
in \(A\) and \(C^*\). Consequently, the incidence functional
is odd. An observable symmetric in both channels preserves its
value. This difference in parity expresses how the same
operation on the structure produces different responses depending
on the reader: areal orientation, symmetric information, and the
mixing vector preserve different information from the same antecedent.


\subsection{Transport of sheet inversion}

Conjugation of the angular charts preserves \(A_{\rm deg}\) and
\(\Delta_{\rm deg}\), and reverses \(h_\theta\). We define
\begin{equation}
 D_\theta=\operatorname{diag}(1,-1,1),\qquad
 c_h=(-5,7,-20)^{\mathsf T},\qquad
 \ell_h=\frac1{3857}(75,176,-150).
 \label{eq:ii-ckm-sheet-data}
\end{equation}
If \(c_A=(2,0,1)^{\mathsf T}\) and
\(c_\Delta=(28,-25/2,-2/3)^{\mathsf T}\), then
\(\ell_hc_A=\ell_hc_\Delta=0\) and \(\ell_hc_h=1\).

\begin{theorem}[Rational reflection and transported metric]
The operator
\begin{equation}
 \mathcal R_{\rm CKM}
 =\mathsf M_{\rm CKM}D_\theta\mathsf M_{\rm CKM}^{-1}
 =I_3-2c_h\ell_h
 \label{eq:ii-ckm-reflection}
\end{equation}
satisfies
\begin{equation}
 \mathcal R_{\rm CKM}^2=I_3,\qquad
 \det\mathcal R_{\rm CKM}=-1,\qquad
 \mathcal R_{\rm CKM}\mathsf M_{\rm CKM}
 =\mathsf M_{\rm CKM}D_\theta.
 \label{eq:ii-ckm-intertwining}
\end{equation}
It is an orthogonal reflection for the positive metric
\begin{equation}
 G_\theta=(\mathsf M_{\rm CKM}^{-1})^{\mathsf T}
                 \mathsf M_{\rm CKM}^{-1},\qquad
 \mathcal R_{\rm CKM}^{\mathsf T}G_\theta\mathcal R_{\rm CKM}=G_\theta.
 \label{eq:ii-ckm-metric}
\end{equation}
\end{theorem}

\begin{proof}
The operator \(P_h=c_h\ell_h\) satisfies \(P_h^2=P_h\), has image
\(\mathbb Rc_h\), and kernel \(\operatorname{span}(c_A,c_\Delta)\).
Thus \((I_3-2P_h)^2=I_3\), with eigenvalues \(1,1,-1\).
Its action on those three columns is that of \(D_\theta\) in the original
basis, proving the conjugation and intertwining.
For \(v\ne0\),
\(v^{\mathsf T}G_\theta v=\|\mathsf M_{\rm CKM}^{-1}v\|^2>0\).
The last identity follows from \(D_\theta^{\mathsf T}D_\theta=I_3\).
\end{proof}

Its complete expression, preserved from the preceding supporting material, is
\begin{equation}
 \mathcal R_{\rm CKM}=
 \begin{pmatrix}
 4607/3857&1760/3857&-1500/3857\\
 -150/551&199/551&300/551\\
 3000/3857&7040/3857&-2143/3857
 \end{pmatrix}.
 \label{eq:ii-ckm-reflection-matrix}
\end{equation}
The phase \(\delta_{\rm deg}=9A_{\rm deg}\) remains fixed under this
reflection. Consequently, sheet transport must not be identified
with CP conjugation of amplitudes, which changes the sign of the complex phase.
The combinatorial parity of a route, its sheet conjugation, and its
physical representation retain the distinct domains in
\cite{cpt,mezcla}.

\subsection{Record parity and CP realization}

To preserve the discrete support of this distinction, let \(\rho\) be a
recorded route of length \(L\), with directions
\(d_t\in\{N,E,S,W\}\) and sheets \(\ell_t\in\{\Sigma,\Pi\}\).
The intermediate states are those produced by TPK transport.
Reflection \(P\) acts by \((i,j)\mapsto(i,-j)\pmod9\) and
permutes \(E,W\); \(C_{\rm hoja}\) permutes \(\Sigma,\Pi\);
\(T_{\rm reg}\) reverses the order of the record and replaces each
direction by its opposite. Inversion is performed on the enriched
record, not on a projection that has forgotten its predecessor.

\begin{proposition}[Invariance of the combinatorial route action]
For \(\lambda\ge0\), the dimensionless functional
\begin{equation}
 \mathcal A_\lambda(\rho)=
 \sum_{t=2}^{L}\left[1-\delta_{d_t,d_{t-1}}
             +\lambda(1-\delta_{\ell_t,\ell_{t-1}})\right]
 \label{eq:ii-discrete-route-action}
\end{equation}
is invariant under \(P\), \(C_{\rm hoja}\), \(T_{\rm reg}\), and
their compositions.
\end{proposition}

\begin{proof}
A global permutation of the alphabet preserves equality between neighbors.
Reversing the order preserves the number of consecutive changes.
Thus each operation separately preserves the changes of direction
and sheet summed in \eqref{eq:ii-discrete-route-action}.
\end{proof}

\(\lambda\) is the declared parameter of this collection of functionals;
\(\mathcal A_\lambda\) is not identified with a dimensional Planck
section. Invariance of the counter also does not imply that every oriented
reader is even. In a mixing realization \(\mathcal V\),
CP compatibility takes the precise form
\begin{equation}
 \mathcal V(CP_{\rm HMT}\rho)
 =D_u\,\overline{\mathcal V(\rho)}\,D_d^\dagger,
 \label{eq:ii-cp-square}
\end{equation}
with the diagonal phase changes and representation spaces
declared. The route counter, rational sheet transport, and
complex conjugation are three distinct maps of the preserved support.

\begin{remark}[CP compatibility of the physical reader]
Equation \eqref{eq:ii-cp-square} is the CP compatibility
condition for the complete physical reader.
The proofs in this section establish invariance of the
record and the identities of the complex reader, but do not
automatically turn \(\mathcal R_{\rm CKM}\) into \(CP_{\rm HMT}\).
This qualification concerns that representation square, not the
existence of the internal route transport already constructed.
\end{remark}

\subsection{Unitary amplitudes, CP phase, and the Jarlskog invariant}

We write \(c_{ab}=\cos\vartheta_{ab}\) and
\(s_{ab}=\sin\vartheta_{ab}\). The complex realization of the reader is
\begingroup
\small
\setlength{\arraycolsep}{3pt}
\begin{equation}
 V_{\rm CKM}=
 \begin{pmatrix}
 c_{12}c_{13}&s_{12}c_{13}&s_{13}e^{-i\delta}\\
 -s_{12}c_{23}-c_{12}s_{23}s_{13}e^{i\delta}&
 c_{12}c_{23}-s_{12}s_{23}s_{13}e^{i\delta}&s_{23}c_{13}\\
 s_{12}s_{23}-c_{12}c_{23}s_{13}e^{i\delta}&
 -c_{12}s_{23}-s_{12}c_{23}s_{13}e^{i\delta}&c_{23}c_{13}
 \end{pmatrix}.
 \label{eq:ii-ckm-unitary}
\end{equation}
\endgroup
The complex structure used is the realization of the oriented
structure already constructed; the unitary parametrization represents the angular
result, rather than determining its coefficients.

\begin{proposition}[Unitarity and phase invariant]
The matrix \eqref{eq:ii-ckm-unitary} is unitary. The quartet
\begin{equation}
 J_{\rm CKM}
 =\operatorname{Im}
   (V_{11}V_{22}\overline{V_{12}}\,\overline{V_{21}})
 =c_{12}c_{23}c_{13}^2s_{12}s_{23}s_{13}\sin\delta
 \label{eq:ii-jarlskog}
\end{equation}
is invariant under \(V\mapsto D_uVD_d\), with \(D_u,D_d\) diagonal
unitary matrices, and changes sign under \(V\mapsto\overline V\).
\end{proposition}

\begin{proof}
The matrix is the product \(R_{23}U_{13}(\delta)R_{12}\), where
\begin{equation}
\begin{aligned}
 R_{12}&=\begin{pmatrix}c_{12}&s_{12}&0\\-s_{12}&c_{12}&0\\0&0&1\end{pmatrix},
 \\
 U_{13}(\delta)&=
 \begin{pmatrix}c_{13}&0&s_{13}e^{-i\delta}\\0&1&0\\
 -s_{13}e^{i\delta}&0&c_{13}\end{pmatrix},
 \\
 R_{23}&=\begin{pmatrix}1&0&0\\0&c_{23}&s_{23}\\0&-s_{23}&c_{23}\end{pmatrix}.
\end{aligned}
 \label{eq:ii-ckm-factorization}
\end{equation}
Each factor is unitary because \(c_{ab}^2+s_{ab}^2=1\).
Substituting the four entries of the quartet, its imaginary part is
\(c_{12}s_{12}c_{13}^2c_{23}s_{23}s_{13}
(c_{12}^2+s_{12}^2)\sin\delta\), giving
\eqref{eq:ii-jarlskog}. The phase factors of each row and each column
cancel with their conjugates. Complex conjugation reverses the imaginary
part.
\end{proof}

A nonzero \(J_{\rm CKM}\) obstructs making the entire matrix real through
those phase changes. Its physical interpretation also retains the corresponding
sector and representation; it is not a scalar correction added
to masses. If the spectra have degeneracies, their change-of-basis
freedoms must be preserved before formulating a CP conclusion.

\subsection{Conjugation of amplitudes and oriented invariant}
\label{sec:ii09-conjugacion-amplitudes}

The complex realization preserves the three mixing angles and
the interference phase as explicit coordinates. Write
\(V(\boldsymbol\vartheta,\delta)
=R_{23}U_{13}(\delta)R_{12}\), with the rotations and phase
defined in the unitary development. The arguments of the
trigonometric functions are their representatives in radians.

\begin{proposition}[Conjugation in the mixing realization]
\label{prop:ii09-conjugacion-unitaria}
For the three declared real angles and the declared real phase,
\[
 V(\boldsymbol\vartheta,-\delta)
       =\overline{V(\boldsymbol\vartheta,\delta)},
 \qquad
 J(\boldsymbol\vartheta,-\delta)
       =-J(\boldsymbol\vartheta,\delta).
\]
The invariant \(J\) also preserves its value under diagonal
phase changes of the charged bases.
\end{proposition}
\begin{proof}
The matrices \(R_{12}\) and \(R_{23}\) have real entries.
In \(U_{13}\), the change \(\delta\mapsto-\delta\) conjugates
the exponential factors. Conjugating the product gives the
first identity. In the quartet
\(V_{11}V_{22}\overline{V_{12}}\overline{V_{21}}\), the
diagonal factors of each row and column cancel. Conjugation
reverses the sign of its imaginary part. Equivalently, the
result follows by substitution into
\[
 J=c_{12}c_{23}c_{13}^{2}s_{12}s_{23}s_{13}\sin\delta.
\]
\end{proof}

In the three-generation sector, the nonzero internal value of
\(J\) expresses the obstruction to realizing all amplitudes
by a real matrix through diagonal phase changes. This
is the charge--parity violation reading of the constructed
matrix. The route reading, in turn, preserves the discrete
operation of conjugation and parity and its representation. Their
full compatibility is expressed by the square
\[
 \mathcal V(CP_{\rm HMT}\rho)
  =D_u\,\overline{\mathcal V(\rho)}\,D_d^{\dagger},
\]
with the representation spaces and phase changes
declared. The conjugation identity proved here fixes
exactly the action that this square must transport.


\Needspace{14\baselineskip}
\begin{samepage}
\begin{proposition}[Mixing and spectrum]
Let \(B_u=\operatorname{diag}(a_1,a_2,a_3)\) and
\(B_d=\operatorname{diag}(b_1,b_2,b_3)\) be the Hermitian mass readers
already produced in their fibers. The transport
\begin{equation}
 B_d^{(u)}=V_{\rm CKM}B_dV_{\rm CKM}^{\dagger}
 \label{eq:ii-ckm-mass-transport}
\end{equation}
preserves the spectrum of \(B_d\), and satisfies
\begin{equation}
 \left|\det[B_u,B_d^{(u)}]\right|
 =2|J_{\rm CKM}|
   \prod_{i<j}|a_i-a_j|\prod_{r<s}|b_r-b_s|.
 \label{eq:ii-ckm-commutator}
\end{equation}
\end{proposition}
\end{samepage}

\begin{proof}
Unitary conjugation preserves the characteristic polynomial. If
\(B=B_d^{(u)}\), the entries of the commutator are
\([B_u,B]_{ij}=(a_i-a_j)B_{ij}\), with zero diagonal. Its determinant is
\[
 2i(a_1-a_2)(a_2-a_3)(a_3-a_1)
     \operatorname{Im}(B_{12}B_{23}B_{31}).
\]
Substituting \(B_{ij}=\sum_r b_rV_{ir}\overline{V_{jr}}\), the imaginary
part is an alternating polynomial of degree three in the \(b_r\).
It is divisible by their three differences: if two \(b_r\) coincide,
\(B=bI+(b'-b)vv^\dagger\), and the off-diagonal product is real.
Expanding a coefficient, using orthogonality of the columns,
gives the quartet in \eqref{eq:ii-jarlskog}, with the sign of the chosen order.
Taking absolute values removes that convention and yields
\eqref{eq:ii-ckm-commutator}, including its vanishing at degeneracies.
\end{proof}

\begin{remark}[Evaluation and comparison]
The decimal CKM coordinates in
Table~\ref{tab:ii-ckm-comparacion} are subsequent evaluations of the
preceding reader. The quantitative comparison section also presents
the angles and moduli of the amplitudes. Those evaluations are not
used to choose its columns or to prove
\eqref{eq:ii-ckm-intertwining} or \eqref{eq:ii-jarlskog}.
This block assigns no new metrological significance to those numbers:
a joint comparison requires its external edition, definition of
observables, and covariance. Marginal deviations do not replace it.
\end{remark}

\subsection{Leptonic record, operators, and spectral frames}

The leptonic sector preserves another domain. The spectral construction of
\eqref{eq:ii-neutrino-compression} brings together four neutral depths of dimensions
\(2,4,6,8\), before the three-dimensional compression:
\begin{equation}
 \mathcal H_\nu^{\rm reg}
 =\bigoplus_{j=1}^{4}\mathcal H_{\nu,G_j},\qquad
 \dim\mathcal H_\nu^{\rm reg}=20,\qquad
 M_\nu=\left.(P_\nu\mathcal M_\nu P_\nu)\right|_{\operatorname{Im}P_\nu},
 \quad\operatorname{rank}P_\nu=3.
 \label{eq:ii-neutrino-compression}
\end{equation}
The four depths are not four flavors or four masses.
\(P_\nu\) is the neutral projector inherited from the record reader, and
\(M_\ell=\left.(P_\ell\mathcal M_\ell P_\ell)\right|_{\operatorname{Im}P_\ell}\),
with \(\operatorname{rank}P_\ell=3\), is the compressed charged operator.
The scale of these operators comes from the action--clock--mass branch;
no PMNS angle determines it.

\begin{definition}[Transport of leptonic frames]
In the declared common three-dimensional realization, let
\(F_\ell,F_\nu:\mathbb C^3\to\mathcal H_3\) be complete orthonormal
spectral frames of the charged and neutral operators. Identification
with \(\mathcal H_3\) is fixed before taking their product. The reader is
\begin{equation}
 U_{\rm PMNS}=F_\ell^\dagger F_\nu:
 \mathbb C^3_\nu\longrightarrow\mathbb C^3_\ell.
 \label{eq:ii-pmns-frames}
\end{equation}
The frames and identification come from the representation and the
record; they are not chosen from an experimental matrix one wishes
to reproduce.
\end{definition}

\begin{theorem}[Unitarity, degeneracies, and naturality]
Under the preceding identification, \(U_{\rm PMNS}\) is unitary.
If the spectra are simple, their frame changes are diagonal phases
and permutations; once spectral order is fixed, the phases remain.
If there are multiplicities \(m_{\ell,\lambda}\) and \(m_{\nu,\eta}\), the
basis-independent object is the double coset
\begin{equation}
 \left(\prod_\lambda U(m_{\ell,\lambda})\right)
 \backslash U(3) /
 \left(\prod_\eta U(m_{\nu,\eta})\right).
 \label{eq:ii-pmns-double-class}
\end{equation}
A common unitary refinement \(W\) intertwining both operators
preserves transport when the frames are prolonged as
\(F'_\ell=WF_\ell\), \(F'_\nu=WF_\nu\).
\end{theorem}

\begin{proof}
Completeness in the same space gives
\(F_\ell F_\ell^\dagger=F_\nu F_\nu^\dagger=I_{\mathcal H_3}\).
Therefore,
\(U_{\rm PMNS}^\dagger U_{\rm PMNS}
=F_\nu^\dagger F_\ell F_\ell^\dagger F_\nu=I_3\), and analogously
in the other order. Basis changes preserving a Hermitian operator
act in blocks within its eigenspaces. Substituting
\(F_\ell D_\ell,F_\nu D_\nu\) yields
\(U\mapsto D_\ell^\dagger UD_\nu\), which determines the double coset.
Finally, \((WF_\ell)^\dagger WF_\nu=F_\ell^\dagger F_\nu\).
\end{proof}

The common-space condition is substantive. For two isometric frames in
a larger ambient space, separate orthonormality yields only a
compression; it does not authorize canceling \(F_\ell F_\ell^\dagger\) as though
it were the ambient identity.

\subsection{Record kernel and rank sufficiency}

The readers \(\Xi_r\) preserve the carry, sheet,
orientation, vacancy, and holonomy components of the record. On the finite set
of admitted routes, one forms
\begin{equation}
 K_{\rm reg}(c,d)=
 \sum_{\substack{r\in\{\mathrm{car},\mathrm{hoj},\mathrm{ori},
                           \mathrm{vac},\mathrm{hol}\}\\ d_r>0}}
 \frac{\langle\Xi_r(c),\Xi_r(d)\rangle}{d_r},
 \qquad
 d_r=\sum_x\|\Xi_r(x)\|^2.
 \label{eq:ii-reg-kernel}
\end{equation}
Normalization comes from the norms of the record itself.
Zero-norm components are omitted; they are not replaced by terms taken
from a target matrix. Each \(\Xi_r\) and each aggregation class must
preserve its route residence and compatibility with refinement.

\begin{proposition}[Positivity and rank of the record]
The matrix \(K_{\rm reg}\) is positive semidefinite. If its
feature vectors are aggregated into three sectors, the resulting Gram
matrix \(G_\nu\) is invertible exactly when those three aggregated
vectors are linearly independent.
\end{proposition}

\begin{proof}
Define
\(v_c=\bigoplus_{r:d_r>0}d_r^{-1/2}\Xi_r(c)\).
Then \(K_{\rm reg}(c,d)=\langle v_c,v_d\rangle\), and for
any coefficients \(z_c\),
\(\sum_{c,d}\overline z_cK_{\rm reg}(c,d)z_d
=\|\sum_c z_cv_c\|^2\ge0\).
The same identity holds for the three aggregated vectors.
The kernel of their Gram matrix consists exactly of their linear
relations, yielding the invertibility criterion.
\end{proof}

That single-sector Gram matrix is distinguished from the \emph{cross} matrix between the
charged classes \(L_a\) and neutral classes \(N_b\), expressed in \cite{mezcla}:
\begin{equation}
 G_{ab}=\frac1{\sqrt{|L_a||N_b|}}
       \sum_{c\in L_a}\sum_{d\in N_b}K_{\rm reg}(c,d).
 \label{eq:ii-cross-gram}
\end{equation}
The nonempty classes come from the internal classifier. The cross matrix
need not be Hermitian, unlike a Gram matrix constructed with
the same three vectors on both sides.

\begin{proposition}[Polar factor and its scope]
If \(G\) is invertible,
\begin{equation}
 U_G=G(G^\dagger G)^{-1/2}
 \label{eq:ii-polar}
\end{equation}
is unitary. If \(G\) is also Hermitian positive definite, then
\(U_G=I_3\). Consequently, the rank of a positive Gram matrix certifies
the sufficiency of the observables, but does not by itself determine
nontrivial mixing between the charged and neutral frames.
\end{proposition}

\begin{proof}
\(G^\dagger G\) is positive definite; its positive square root is invertible, and
\[
 U_G^\dagger U_G=(G^\dagger G)^{-1/2}(G^\dagger G)
 (G^\dagger G)^{-1/2}=I_3.
\]
By finite dimensionality, also
\(U_GU_G^\dagger=I_3\). If \(G>0\), the positive square root of
\(G^\dagger G=G^2\) is \(G\), and \(U_G=I_3\).
\end{proof}

\begin{remark}[Polar factorization and spectral frames]
Equations \eqref{eq:ii-reg-kernel} and \eqref{eq:ii-cross-gram}
preserve, respectively, the record's rank criterion and
the cross matrix of the charged and neutral sectors. This exposition
separates their functions to avoid turning a positive Gram matrix into an artificial
angle selector. Identifying \(U_G\) with
\eqref{eq:ii-pmns-frames} requires compatibility of the cross kernel
with the declared spectral frames. The identity of a polar factor
or mere unitarity does not replace that identification. No new
numerical PMNS matrix is specified here.
\end{remark}

\subsection{Check of the rank-one phase}

Let \(b_K\) be the source's normalized incidence witness and
\(Q_K=|b_K\rangle\langle b_K|\). The operator
\begin{equation}
 \Phi_K=I+(e^{i\pi_{\rm HMT}/6}-1)Q_K
 \label{eq:ii-rank-one-phase}
\end{equation}
is unitary in its ambient space. For isometric frames, consider
\(U_\Phi=F_\ell^\dagger\Phi_KF_\nu\).

\begin{proposition}[Exact criterion for a unitary compression]
The operator \(U_\Phi\) is a contraction. It is unitary exactly
when
\begin{equation}
 \Phi_K(\operatorname{Im}F_\nu)=\operatorname{Im}F_\ell.
 \label{eq:ii-phase-compatibility}
\end{equation}
\end{proposition}

\begin{proof}
Let \(P_\ell=F_\ell F_\ell^\dagger\) and
\(P_\nu=F_\nu F_\nu^\dagger\). Then
\(U_\Phi^\dagger U_\Phi
=F_\nu^\dagger\Phi_K^\dagger P_\ell\Phi_KF_\nu\le I_3\).
Equality is equivalent to
\((I-P_\ell)\Phi_KF_\nu=0\), that is, to inclusion of the
left-hand image in \eqref{eq:ii-phase-compatibility} in the right-hand one.
Both have dimension three, so the inclusion is equality.
\end{proof}

The source \cite{mezcla} preserves this minimal ansatz as a negative
check on physical identification. Its ambient unitarity, and even
satisfaction of \eqref{eq:ii-phase-compatibility}, do not replace the
complete record or the frames of \eqref{eq:ii-pmns-frames}. The
check is not removed when angular propagation is assembled in article II.

\subsection{Leptonic phase and the Dirac--Majorana criterion}

Phase changes of bases must be distinguished from the self-conjugation
data of the neutral sector. For any unitary leptonic
transport, one can form the quartet
\begin{equation}
 J_\ell=\operatorname{Im}
 \bigl(U_{e1}U_{\mu2}\overline{U_{e2}}\,
                        \overline{U_{\mu1}}\bigr).
 \label{eq:ii-leptonic-quartet}
\end{equation}
The phase cancellation proved for \eqref{eq:ii-jarlskog}
applies equally here. The quartet does not fix all the data of the
neutral realization or determine by itself whether it is of Dirac or Majorana type.

The Majorana realization is specified by the
antiunitary self-conjugation \(\mathcal C_\nu\) compatible with mass and
dynamics:
\begin{equation}
 \mathcal C_\nu^2=I,\qquad
 \mathcal C_\nu M_\nu=M_\nu\mathcal C_\nu,\qquad
 \mathcal C_\nu H_\nu=H_\nu\mathcal C_\nu,
 \label{eq:ii-majorana-compatibility}
\end{equation}
with domains preserved. The representation is realized in the corresponding
fixed real form. The Dirac alternative also preserves a charge
\(Q\) distinguishing the mode from its conjugate:
\begin{equation}
 \mathcal C_\nu Q\mathcal C_\nu^{-1}=-Q,\qquad Q\psi\ne0.
 \label{eq:ii-dirac-charge}
\end{equation}
These statements preserve their representational data. Neutrality,
the involution of dodecaphasic words, and a topological Ising sector
do not separately decide the alternative.

\begin{remark}[Basis freedom and Majorana phases]
The double coset \eqref{eq:ii-pmns-double-class} describes the basis
freedom of the Hermitian pair. If a Majorana real form is additionally fixed,
admissible phase changes must preserve that additional datum;
not all the freedoms of a Dirac representation can be transferred
without qualification. The present block preserves the self-conjugation criterion and
assigns no values to Majorana phases or neutrino masses. Such values
are not obtained from neutrality or from a rank-one polarization.
\end{remark}

\Needspace{8\baselineskip}
\subsection{Closure of angular transport}

The quark construction expresses the angular coordinates first, then
the complex amplitudes, and finally the phase invariants and
operator transport. The leptonic construction expresses record
operators and spectral frames before the change of basis. Both routes
preserve their residence in the common enriched state and in the discrete
structure of the continuum; neither introduces an experimental table as
a selector of coefficients or states.

The checks in this section are local and explicit: rational CKM inversion and
intertwining; spectral preservation by conjugation;
completeness and common space of the PMNS frames; rank of the record;
compatibility of the cross matrix; and preservation of the self-conjugation
data. Angular propagation is thus incorporated into the theory of
action, area, and information without making mixing a second mass
law or turning the Barbero reader into its cause.
